\documentclass[11pt]{article}

\usepackage[margin=25mm]{geometry}

\usepackage{amsmath,amssymb,fontspec,unicode-math}

\IfFontExistsTF{DejaVu Serif}{\setmainfont{DejaVu Serif}}{\setmainfont{Latin Modern Roman}}

\IfFontExistsTF{Latin Modern Math}{\setmathfont{Latin Modern Math}}{\setmathfont{DejaVu Math TeX Gyre}}

\usepackage{xurl,hyperref}

\setlength{\parindent}{0pt}\setlength{\parskip}{6pt}

\begin{document}

\title{Revisiting Semi-Supervised Learning in the Era of Foundation Models}

\author{ELVIA preliminary AI analysis}\date{21 October 2026}\maketitle

Preliminary assessment: Minor.\par

Independent scientific review remains pending. Author adoption has not been requested.

The default release selects teachers using squared average ranks and transforms vote counts by softmax, differing from the written average-rank and mean-label operations. Exact inputs change selection and distillation gradients; historical benchmark selections remain unresolved.

\section*{Review scope}

Claim-focused reading of the original paper and relevant author-released passages; direct mathematical or static-source reasoning. No author program, earlier verifier, model or benchmark was executed during this check.

\section*{Problem 1: Mean squared rank can select a different teacher}

Source: Section 3.3, PDF p.5; Appendix C.1; tune\_print.py and gen\_config\_pet.py rank branches

Assessment: supported; impact: minor

\subsection*{Source claim (paraphrase)}

The paper minimizes average validator rank; the default and teacher-generation branches square ranks first.

\subsection*{Analysis}

Across seven tie-free validators, let A ranks be \(1,1,1,1,1,3,3\), B ranks \(2,2,2,2,2,1,1\), C ranks \(3,3,3,3,3,2,2\). Ordinary means are \(11/7,12/7,19/7\), uniquely choosing A. Squared means are \(23/7,22/7,53/7\), choosing B. Each validator is a valid strict permutation. This is an operator-input example, without asserting the exact joint scores arise in the reported images. Squared mean equals squared mean rank plus across-validator variance, adding a different preference.

\subsection*{Finding and scope}

The two selection rules are inequivalent. A non-squared reporting option exists, but does not change the separately squared teacher-generation branches; actual published choices remain unknown.

\subsection*{Proposed correction}

Follow the stated mean-rank operation in every relevant branch, or explicitly define and attribute squared-rank selection with its dispersion sensitivity.

A proposed correction is retained where stated; independent verification and author adoption remain unestablished.

\section*{Problem 2: Softmax of counts is not a mean one-hot label}

Source: Algorithm 1, PDF p.7; PETHook.init\_pl voting branch and distillation target path

Assessment: supported; impact: minor

\subsection*{Source claim (paraphrase)}

The stated target is \(v/K\), whereas the release passes \(\operatorname{softmax}(v)\) into the soft distillation loss.

\subsection*{Analysis}

Four teachers voting for classes \(0,0,1,2\) give counts \((2,1,1,0)\). Mean labels are \((0.5,0.25,0.25,0)\); the release gives approximately \[(0.534447,0.196612,0.196612,0.072329)\]. At zero student logits, the KL gradient is student probability minus target: the written target gives \((-0.25,0,0,0.25)\), while the release gives \[(-0.284447,0.053388,0.053388,0.177671)\]. Argmax labels agree, but soft training differs. Duplicating every teacher preserves vote proportions while sharpening the released distribution.

\subsection*{Finding and scope}

The release retains scale-insensitive teacher argmax voting but uses a different soft target and optimization path. No actual accuracy change is established.

\subsection*{Proposed correction}

Use \(v/K\) to restore Mean Label, or explicitly specify the count-dependent softmax target and distinguish that design from average labels.

A proposed correction is retained where stated; independent verification and author adoption remain unestablished.

\section*{Limits}

Preliminary AI assessment; independent human scientific adjudication remains pending.\par

Analytic witnesses and static release checks do not remeasure the published benchmark results.\par

\section*{Source identity}

NeurIPS 2025 official proceedings PDF (retained local copy)

\url{https://proceedings.nips.cc/paper_files/paper/2025/file/558159930b584e7f137a1dcbd380d9fd-Paper-Conference.pdf}

Reviewed PDF SHA-256: \nolinkurl{6c5cad993bc1e3dd0e0ddde8174f036e6a371049ec8eafb3c7c8997210c525cd}

The SHA-256 identifies the exact retained PDF inspected in this rapid source/argument check. Original supplied report bytes are retained locally. The original PDF retrieval timestamp is taken from matching retained metadata.

\end{document}
